\documentclass[11pt]{article}
\usepackage[margin=1.1in]{geometry}
\usepackage{amsmath,amssymb,amsthm,mathtools,booktabs,array,tabularx,longtable}
\usepackage[dvipsnames]{xcolor}
\usepackage[colorlinks=true,linkcolor=MidnightBlue,citecolor=MidnightBlue,urlcolor=MidnightBlue]{hyperref}
\usepackage{fancyhdr}
\pagestyle{fancy}\fancyhf{}\renewcommand{\headrulewidth}{0pt}
\fancyfoot[R]{\thepage}
\fancyfoot[L]{\scriptsize\color{gray}Generated by the Transport Lab from approved parts of its store; not a deposited paper. Papers of record: \href{https://zenodo.org/communities/transport-papers}{Zenodo, Transport Papers}}
\fancypagestyle{plain}{\fancyhf{}\fancyfoot[R]{\thepage}\fancyfoot[L]{\scriptsize\color{gray}Generated by the Transport Lab from approved parts of its store; not a deposited paper. Papers of record: \href{https://zenodo.org/communities/transport-papers}{Zenodo, Transport Papers}}\renewcommand{\headrulewidth}{0pt}}
\newtheorem{theorem}{Theorem}[section]
\newtheorem{proposition}[theorem]{Proposition}
\newtheorem{lemma}[theorem]{Lemma}
\newtheorem{corollary}[theorem]{Corollary}
\newtheorem{conjecture}[theorem]{Conjecture}
\newtheorem{refutation}[theorem]{Refutation}
\newtheorem{prediction}[theorem]{Prediction}
\newtheorem{principle}[theorem]{Principle}
\newtheorem{observation}[theorem]{Observation}
\newtheorem{axiom}[theorem]{Axiom}
\theoremstyle{definition}
\newtheorem{definition}[theorem]{Definition}
\newtheorem{fact}[theorem]{Fact}
\newtheorem{assembled}[theorem]{Assembly}
\newtheorem{postulate}[theorem]{Postulate}
\newtheorem{premise}[theorem]{Premise}
\newtheorem{hypothesis}[theorem]{Hypothesis}
\newtheorem{remark}[theorem]{Remark}
\newtheorem{example}[theorem]{Example}
\theoremstyle{plain}
\setlength{\marginparwidth}{0.8in}
\newcommand{\changedmark}[2]{\marginpar{\raggedright\scriptsize\color{#1}#2}}
\providecommand{\Aa}{}\renewcommand{\Aa}{\mathcal{A}}
\providecommand{\Ad}{}\renewcommand{\Ad}{\widetilde{\mathrm{Ad}}}
\providecommand{\Aev}{}\renewcommand{\Aev}{\mathcal{A}^{+}}
\providecommand{\Aod}{}\renewcommand{\Aod}{\mathcal{A}^{-}}
\providecommand{\B}{}\renewcommand{\B}{\mathrm{B}}
\providecommand{\C}{}\renewcommand{\C}{\mathbb{C}}
\providecommand{\CAR}{}\renewcommand{\CAR}{\mathrm{CAR}}
\providecommand{\CS}{}\renewcommand{\CS}{{\mathbb{S}^{+}}}
\providecommand{\CSm}{}\renewcommand{\CSm}{{\mathbb{S}^{-}}}
\providecommand{\Cl}{}\renewcommand{\Cl}{\mathrm{Cl}}
\providecommand{\Con}{}\renewcommand{\Con}{\mathcal{C}}
\providecommand{\Dl}{}\renewcommand{\Dl}{\Delta}
\providecommand{\EN}{}\renewcommand{\EN}{E_{N}}
\providecommand{\End}{}\renewcommand{\End}{\mathrm{End}}
\providecommand{\F}{}\renewcommand{\F}{\mathbb{F}}
\providecommand{\G}{}\renewcommand{\G}{\Gamma}
\providecommand{\HH}{}\renewcommand{\HH}{\mathbb{H}}
\providecommand{\IS}{}\renewcommand{\IS}{I\sigma_3}
\providecommand{\J}{}\renewcommand{\J}{\mathbf{J}}
\providecommand{\LH}{}\renewcommand{\LH}{L_{H}}
\providecommand{\Mat}{}\renewcommand{\Mat}{\mathrm{M}}
\providecommand{\Mo}{}\renewcommand{\Mo}{M^{\circ}}
\providecommand{\Mp}{}\renewcommand{\Mp}{M^{+}}
\providecommand{\Msun}{}\renewcommand{\Msun}{M_{\odot}}
\providecommand{\Mt}{}\renewcommand{\Mt}{\widetilde{M}}
\providecommand{\Ncone}{}\renewcommand{\Ncone}{\mathcal{N}}
\providecommand{\Neff}{}\renewcommand{\Neff}{N_{\mathrm{eff}}}
\providecommand{\PCS}{}\renewcommand{\PCS}{\mathbb{P}{\mathbb{S}^{+}}}
\providecommand{\PT}{}\renewcommand{\PT}{\mathbb{PT}}
\providecommand{\Pin}{}\renewcommand{\Pin}{\mathrm{Pin}}
\providecommand{\Q}{}\renewcommand{\Q}{\mathbb{Q}}
\providecommand{\R}{}\renewcommand{\R}{\mathbb{R}}
\providecommand{\Ric}{}\renewcommand{\Ric}{\mathrm{Ric}}
\providecommand{\SBH}{}\renewcommand{\SBH}{S_{\mathrm{BH}}}
\providecommand{\Sc}{}\renewcommand{\Sc}{\mathbb{S}}
\providecommand{\Sch}{}\renewcommand{\Sch}{\mathrm{P}}
\providecommand{\Sig}{}\renewcommand{\Sig}{\Sigma}
\providecommand{\Som}{}\renewcommand{\Som}{S_\omega}
\providecommand{\Sp}{}\renewcommand{\Sp}{\mathbb{S}}
\providecommand{\Sph}{}\renewcommand{\Sph}{\mathrm{S}}
\providecommand{\Spi}{}\renewcommand{\Spi}{S_\pi}
\providecommand{\Spin}{}\renewcommand{\Spin}{\mathrm{Spin}}
\providecommand{\Sth}{}\renewcommand{\Sth}{S_{\mathrm{th}}}
\providecommand{\T}{}\renewcommand{\T}{\mathbb{T}}
\providecommand{\TH}{}\renewcommand{\TH}{T_{\mathrm{H}}}
\providecommand{\Tr}{}\renewcommand{\Tr}{\operatorname{Tr}}
\providecommand{\UHF}{}\renewcommand{\UHF}{\mathrm{UHF}}
\providecommand{\Z}{}\renewcommand{\Z}{\mathbb{Z}}
\providecommand{\Zc}{}\renewcommand{\Zc}{\mathrm{Z}}
\providecommand{\cmg}{}\renewcommand{\cmg}{\mathrm{cm^{2}\,g^{-1}}}
\providecommand{\dS}{}\renewcommand{\dS}{\mathrm{dS}}
\providecommand{\dsc}{}\renewcommand{\dsc}{\mathrm{dsc}}
\providecommand{\dtPG}{}\renewcommand{\dtPG}{dt}
\providecommand{\dtS}{}\renewcommand{\dtS}{dt_{\mathrm{S}}}
\providecommand{\ellv}{}\renewcommand{\ellv}{\ell}
\providecommand{\er}{}\renewcommand{\er}{e_r}
\providecommand{\et}{}\renewcommand{\et}{e_t}
\providecommand{\ev}{}\renewcommand{\ev}{\mathrm{ev}}
\providecommand{\gl}{}\renewcommand{\gl}{\mathfrak{gl}}
\providecommand{\hatV}{}\renewcommand{\hatV}{\widehat{V}}
\providecommand{\hb}{}\renewcommand{\hb}{\bar{h}}
\providecommand{\hf}{}\renewcommand{\hf}{\hat{f}}
\providecommand{\id}{}\renewcommand{\id}{\mathrm{id}}
\providecommand{\kpc}{}\renewcommand{\kpc}{\mathrm{kpc}}
\providecommand{\lP}{}\renewcommand{\lP}{\ell_{P}}
\providecommand{\mD}{}\renewcommand{\mD}{\mathbf{m}_{D}}
\providecommand{\mb}{}\renewcommand{\mb}{\bar m}
\providecommand{\mnu}{}\renewcommand{\mnu}{\Sigma m_{\nu}}
\providecommand{\nInf}{}\renewcommand{\nInf}{n_\infty}
\providecommand{\nO}{}\renewcommand{\nO}{n_0}
\providecommand{\ninf}{}\renewcommand{\ninf}{n_{\infty}}
\providecommand{\nz}{}\renewcommand{\nz}{n_{0}}
\providecommand{\proj}{}\renewcommand{\proj}{\operatorname{proj}}
\providecommand{\pv}{}\renewcommand{\pv}{\pi}
\providecommand{\rad}{}\renewcommand{\rad}{\mathrm{rad}}
\providecommand{\sgn}{}\renewcommand{\sgn}{\operatorname{sgn}}
\providecommand{\shA}{}\renewcommand{\shA}{\mathcal{A}}
\providecommand{\shC}{}\renewcommand{\shC}{\mathcal{C}}
\providecommand{\shE}{}\renewcommand{\shE}{\mathcal{E}}
\providecommand{\shJ}{}\renewcommand{\shJ}{\mathcal{J}}
\providecommand{\shK}{}\renewcommand{\shK}{\mathcal{K}}
\providecommand{\shL}{}\renewcommand{\shL}{L}
\providecommand{\shW}{}\renewcommand{\shW}{\mathcal{W}}
\providecommand{\sll}{}\renewcommand{\sll}{\mathfrak{sl}}
\providecommand{\so}{}\renewcommand{\so}{\mathfrak{so}}
\providecommand{\su}{}\renewcommand{\su}{\mathfrak{su}}
\providecommand{\supp}{}\renewcommand{\supp}{\operatorname{supp}}
\providecommand{\tr}{}\renewcommand{\tr}{\operatorname{tr}}
\providecommand{\uu}{}\renewcommand{\uu}{\mathfrak{u}}
\providecommand{\br}{}\renewcommand{\br}[1]{\langle #1\rangle}
\providecommand{\bra}{}\renewcommand{\bra}[1]{\langle #1|}
\providecommand{\file}{}\renewcommand{\file}[1]{\texttt{#1}}
\providecommand{\hp}{}\renewcommand{\hp}[2]{( #1 \,|\, #2 )}
\providecommand{\ip}{}\renewcommand{\ip}[2]{\langle #1, #2\rangle}
\providecommand{\ket}{}\renewcommand{\ket}[1]{|#1\rangle}
\providecommand{\rel}{}\renewcommand{\rel}[2]{D\bigl(#1 \,\|\, #2\bigr)}
\providecommand{\rev}{}\renewcommand{\rev}[1]{\widetilde{#1}}
\providecommand{\slashed}{}\renewcommand{\slashed}[1]{{#1}\mkern-10mu/}
\title{\bfseries Degenerate metrics and walls}
\author{Leon Butler\\ \normalsize\href{mailto:lbutler.math@gmail.com}{\texttt{lbutler.math@gmail.com}}\\ \normalsize ORCID: \href{https://orcid.org/0009-0006-4449-5617}{0009-0006-4449-5617}}
\date{Generated by the Transport Lab, 2 November 2026}
\begin{document}
\maketitle
\begin{abstract}
\normalsize
Geometry can degenerate: along a wall, one direction can lose its length. This trail follows what then happens to the Clifford algebra attached to each point, assuming the degeneration is tame, the mildest controlled kind.

The algebras keep the same size everywhere, so degeneration shows up in their structure rather than in a count of dimensions. On the wall the dying direction $k$ generates a square-zero ideal of the algebra.

Adjoining a hyperbolic plane coupled to the dying direction repairs the geometry on one side, the wall included. The degeneracy is not destroyed, though: it reappears at a shadow wall further away, and for a compact wall with good couplings the process repeats forever.
\end{abstract}
\newpage
\noindent \textbf{Authorship and disclosure. }The founding idea of the series is the author's: that division by zero in a Clifford algebra should be read not as a prohibition but as a transport into a higher-dimensional algebra, and that the transported division need not be accompanied by an inverse operation. The direction of investigation and the critical review of the results are likewise the author's. The mathematical development and the drafting of the text were carried out in extended collaboration with Claude Fable 5 (Anthropic), an AI system.

\noindent\fbox{\parbox{\dimexpr\linewidth-2\fboxsep-2\fboxrule}{\textbf{This is not a deposited paper.} It was generated by the Transport Lab on 2 November 2026 and assembled from approved parts of its store. It follows results across several papers of the series; each is stated with its proof and credited to the paper that first proved it. The official versions of those papers are on Zenodo, in the Transport Papers community, \url{https://zenodo.org/communities/transport-papers}.}}
\newpage

\section*{Introduction}
After the standing assumption, the walk goes from the algebra on a degenerate wall, to its repair, to what the repair reveals about the wall, and finally to the shadow walls the repair creates.

\section{What we assume}

\begin{postulate}[Tameness]\label{c:sheaf.tameness}
This is the standing assumption on how a geometry is allowed to degenerate. The places where it degenerates must form a smooth closed wall, exactly one direction must lose its length at each point of the wall, and these directions must vary smoothly along it. In addition, the length must vanish at a linear rate as the wall is reached, not to higher order.
\end{postulate}

\section{Degenerate metrics and walls}

Each step either describes the wall, repairs it, or follows the consequences of the repair.

We begin before any wall is singled out.

Gluing together the Clifford algebras of a cometric $G$ at all points gives a sheaf $\shA$ that is locally free of constant rank $2^n$, degenerate or not, with structure functions polynomial in $G$. So degeneration cannot be detected by counting dimensions; it has to show up in the structure of the algebra.

\begin{proposition}[Local freeness at constant rank]\label{c:sheaf.local-freeness-constant-rank}
Gluing together the Clifford algebras at all points of a space gives a smoothly varying family of algebras. This result shows that the family has the same size everywhere, even where the geometry degenerates. So degeneration cannot be detected by counting dimensions; it has to show up in the structure of the algebra instead.
\emph{First proved in} T4, Prop 2.8 (\href{https://doi.org/10.5281/zenodo.22721669}{doi:10.5281/zenodo.22721669}).
\end{proposition}
\noindent\emph{Rests on:} \ref{c:sheaf.blade-basis-degenerate-forms-included}.

This step finds that structure, for a tame corank-one degeneration: one direction losing its length, transversally, along a smooth wall $\Sig$.

At a wall point $x$ the radical line, spanned by $k$, generates a two-sided ideal with \[ J(x) = A(x)\,k, \qquad J(x)\^{}2 = 0, \qquad \dim\_\R J(x) = 2\^{}{n-1} \] The quotient is the Clifford algebra of the induced nondegenerate form, and these ideals assemble into a sheaf supported exactly on the wall.

\begin{theorem}[The degeneration sheaf]\label{c:sheaf.degeneration-sheaf}
Where the geometry degenerates along a wall, assuming it does so in the controlled way the paper calls tame (one direction losing its length, transversally, along a smooth wall), the algebra attached to each point changes character. Away from the wall it is an ordinary Clifford algebra. On the wall, the direction that has lost its length generates a part of the algebra whose products all vanish, an infinitesimal part of exactly half the size, and removing it leaves a smaller ordinary Clifford algebra. These parts fit together smoothly along the wall into a well-behaved algebraic object living only on the wall.
\emph{First proved in} T4, Thm 3.1 (\href{https://doi.org/10.5281/zenodo.22721669}{doi:10.5281/zenodo.22721669}).
\end{theorem}
\noindent\emph{Rests on:} \ref{c:sheaf.blade-basis-degenerate-forms-included}, \ref{c:sheaf.imported-classification-values}, \ref{c:std-algebra.vector-squares}, \ref{c:primer.commutation-rule}, \ref{c:std-algebra.existence-uniqueness}, \ref{c:std-algebra.radical-ideal-and-quotient}, \ref{c:sheaf.tameness}.

The wall can also be approached from outside rather than examined in place.

Along a transversal with parameter $t$, rescaling the dying direction as $k \mapsto k/\sqrt{t}$ gives a nondegenerate limit form \[ \hat G = G\_0|\_{V\_1} \perp \dot G(k,k) \] in which the dying direction re-enters with its rate of degeneration as its coefficient. Its Clifford algebra, the exceptional algebra, is canonically independent of the choices made to build it.

\begin{proposition}[The exceptional algebra: deformation to the normal cone]\label{c:sheaf.exceptional-algebra-deformation-normal-cone}
Assume the wall is tame, so that at the wall one direction of the geometry loses its length at a nonzero rate. If that direction is stretched at exactly the right rate as the wall is approached, the geometry has a well-behaved, non-degenerate limit. The algebra of this limit, the exceptional algebra, depends only on the wall point and not on the choices made to build it.
\emph{First proved in} T4, Prop 6.1 (\href{https://doi.org/10.5281/zenodo.22721669}{doi:10.5281/zenodo.22721669}).
\end{proposition}
\noindent\emph{Rests on:} \ref{c:sheaf.blade-basis-degenerate-forms-included}, \ref{c:sheaf.degeneration-sheaf}, \ref{c:sheaf.tameness}.

The repair begins by adjoining a hyperbolic plane $(\nO, \nInf)$ and coupling it to the dying direction with constants $\mu$ and $\nu$.

On the wall the surviving directions split off, and the coupled block has Gram matrix \[ \begin{pmatrix} 0 \& \nu \& \mu \\ \nu \& 0 \& -1 \\ \mu \& -1 \& 0 \end{pmatrix}, \qquad \det = -2\mu\nu \] which is nondegenerate for every nonzero $\mu, \nu$. The sign of $\mu\nu$ decides the signature: $(p{+}1, q{+}1)$ when $\mu\nu > 0$, $(p, q{+}2)$ when $\mu\nu < 0$.

\begin{lemma}[The wall Gram]\label{c:sheaf.wall-gram}
On the wall, the enlarged geometry splits into two independent pieces: the ordinary directions that survive the degeneration, and a three-dimensional block made of the dying direction and the two added ones. A direct calculation shows that this block is never degenerate, whatever non-zero coupling strengths are used. The sign of the product of the two coupling constants decides how many directions of each sign the repaired geometry has.
\emph{First proved in} T4, Lem 4.2 (\href{https://doi.org/10.5281/zenodo.22721669}{doi:10.5281/zenodo.22721669}).
\end{lemma}
\noindent\emph{Rests on:} \ref{c:sheaf.degeneration-sheaf}, \ref{c:std-algebra.sylvester-inertia}.

The block shows that the repair works on the wall; this step extends it to the whole physical side.

Assuming tameness, $p \geq 1$ on the physical side, and a trivial radical line bundle, the coupling gives an ambient form of constant signature $(p{+}1, q{+}1)$ on the whole physical side, wall included. In the physical case every fibre of the completed sheaf is \[ \Cl(2,4) \cong \Mat\_4(\HH) \]

\begin{theorem}[One-sided completion]\label{c:sheaf.one-sided-completion}
Assume a geometry degenerates along a wall in the tame way, with at least one positive direction on the physical side and with the dying direction's line consistently orientable along the wall. Then it can be repaired on that side by adding two extra dimensions and coupling them to the direction that is dying. This result shows that the coupling can be chosen so that the enlarged geometry is non-degenerate, with the same signature everywhere on the physical side, the wall itself included. The resulting family of algebras is then uniform: in the physical case every member is the algebra of four-by-four quaternion matrices.
\emph{First proved in} T4, Thm 4.4 (\href{https://doi.org/10.5281/zenodo.22721669}{doi:10.5281/zenodo.22721669}).
\end{theorem}
\noindent\emph{Rests on:} \ref{c:sheaf.determinant-identity-blow-up-expansion}, \ref{c:sheaf.local-freeness-constant-rank}, \ref{c:sheaf.wall-gram}, \ref{c:sheaf.tameness}.

Next comes what the repair does to the original degenerate algebra.

It sits inside the completed algebra unchanged, but on the wall the radical covector $k$ acquires a partner $k^*$ built from the coupling: \[ k k\^{}* + k\^{}* k = 1, \qquad k\^{}2 = (k\^{}*)\^{}2 = 0 \] The two form a fermionic ladder pair, and this pairing exists exactly over the wall.

\begin{corollary}[Localised transport]\label{c:sheaf.localised-transport}
When the degenerate algebra is placed inside the larger completed algebra, it sits there unchanged. What changes is its surroundings: on the wall, a direction of zero length that was a dead end in the small algebra finds a partner in the large one, and together they behave like the creation and annihilation operators of a single fermion. This pairing happens only on the wall, so the fermion-like structure lives exactly where the geometry is singular.
\emph{First proved in} T4, Cor 4.5 (\href{https://doi.org/10.5281/zenodo.22721669}{doi:10.5281/zenodo.22721669}).
\end{corollary}
\noindent\emph{Rests on:} \ref{c:sheaf.blade-basis-degenerate-forms-included}, \ref{c:sheaf.degeneration-sheaf}, \ref{c:sheaf.determinant-identity-blow-up-expansion}, \ref{c:sheaf.resolution-ideal-theoretically}, \ref{c:sheaf.wall-gram}, \ref{c:sheaf.one-sided-completion}, \ref{c:std-algebra.existence-uniqueness}, \ref{c:std-algebra.sylvester-inertia}.

The repair also says something about the geometry of the wall itself.

Near a tame wall the physical metric, the inverse cometric, has a canonical simple pole whose residue is the inverse of the transversality datum. That residue pairs the conormal line with the dying line twice over, \[ N\^{}*\Sig \cong \shL\^{}{\otimes 2} \] so the completion's partner line is a square root of the normal bundle.

\begin{theorem}[The metric's residue, and the square-root twist]\label{c:sheaf.metrics-residue-square-root-twist}
Near a degenerate wall of the well-behaved kind the series assumes (a tame wall, where the degeneration is transverse), the ordinary metric blows up like one over the distance to the wall. The coefficient of this blow-up turns out to be the inverse of the canonical rate at which the geometry degenerates. It shows that the direction across the wall is, in a precise sense, the square of the direction that dies, so the extra direction added in the repair is a square root of the direction across the wall.
\emph{First proved in} T4, Thm 6.2 (\href{https://doi.org/10.5281/zenodo.22721669}{doi:10.5281/zenodo.22721669}).
\end{theorem}
\noindent\emph{Rests on:} \ref{c:sheaf.invariant-form-t3}, \ref{c:sheaf.localised-transport}, \ref{c:sheaf.tameness}.

That square-root relation has a topological consequence.

A real line bundle may twist like a band with a half twist, but two twists cancel, so the square of any line bundle is untwisted. Hence \[ w\_1(N\^{}*\Sig) = 2\,w\_1(\shL) = 0 \] and every tame wall is two-sided, so a physical side always exists.

\begin{corollary}[Tame walls are two-sided]\label{c:sheaf.tame-walls-are-two-sided}
Some surfaces, like a M\"obius band, have only one side. This result shows that, if a geometry degenerates along a wall in the tame way, as the tameness hypothesis requires, the wall can never be one-sided. So it always makes sense to speak of the physical side of the wall and the other side.
\emph{First proved in} T4, Cor 7.1 (\href{https://doi.org/10.5281/zenodo.22721669}{doi:10.5281/zenodo.22721669}).
\end{corollary}
\noindent\emph{Rests on:} \ref{c:sheaf.imported-real-line-bundles}, \ref{c:sheaf.metrics-residue-square-root-twist}, \ref{c:sheaf.tameness}.

Having two sides matters, because the repair is one-sided.

With $\mu\nu > 0$ the ambient form degenerates again exactly where \[ F = 1 + 2\mu\nu\beta\^{}2 s \] vanishes, a shadow wall on the far side across which the signature passes from $(p{+}1,q{+}1)$ to $(p,q{+}2)$. Reversing the sign only moves the shadow wall to the other side; it never removes it.

\begin{proposition}[Conservation of degeneracy, located]\label{c:sheaf.conservation-degeneracy-located}
The one-sided repair of a degenerating geometry works on one side of the wall, but the degeneracy reappears as a new shadow wall further away. This result locates that shadow wall exactly, as the zero set of an explicit function, and says how the signature changes across it. Flipping the sign of the coupling moves the shadow wall to the other side but never removes it.
\emph{First proved in} T4, Prop 5.2 (\href{https://doi.org/10.5281/zenodo.22721669}{doi:10.5281/zenodo.22721669}).
\end{proposition}
\noindent\emph{Rests on:} \ref{c:sheaf.determinant-identity-blow-up-expansion}, \ref{c:sheaf.wall-gram}.

A shadow wall invites the same repair again.

For a compact tame wall and a gentle enough coupling, the shadow wall is just as tame as the original. Repeating the repair gives an infinite sequence of walls whose fibres ascend the Witt tower, \[ \Cl(p,q) \to \Cl(p{+}1,q{+}1) \to \Cl(p{+}2,q{+}2) \to \cdots \] one rung per wall. The degeneracy is pushed further from the physical region at every step but never destroyed.

\begin{corollary}[The infinite wall sequence]\label{c:sheaf.infinite-wall-sequence}
Provided the wall is compact and each repair is done gently enough (a good coupling at every stage), repairing a degenerate wall on one side creates a new wall a little further away, and that wall can be repaired in the same way. Repeating this forever gives an infinite chain of walls, each carrying an algebra built on two more dimensions than the last. The degeneracy is never destroyed, only pushed further from the physical region at each step.
\emph{First proved in} T4, Cor 5.5 (\href{https://doi.org/10.5281/zenodo.22721669}{doi:10.5281/zenodo.22721669}).
\end{corollary}
\noindent\emph{Rests on:} \ref{c:sheaf.one-sided-completion}, \ref{c:sheaf.tame-wall-induction}.

\section{Results used from elsewhere}
The proofs above use the following results of the series, stated briefly; each is proved in the paper named.

\begin{fact}[The commutation rule]\label{c:primer.commutation-rule}
If $x$ lies in the subalgebra generated by vectors orthogonal to $v$, then $v x = \hat{x} v$ and $x v = v \hat{x}$.
\emph{Source:} T0, Lem 9.1 (\href{https://doi.org/10.5281/zenodo.22714501}{doi:10.5281/zenodo.22714501}).
\end{fact}

\begin{fact}[Blade basis, degenerate forms included]\label{c:sheaf.blade-basis-degenerate-forms-included}
For any symmetric bilinear form, degenerate or not, and any basis, the $2^n$ blades form a basis of the Clifford algebra, with structure constants polynomial in the entries of the form.
\emph{Source:} T4, Lem 2.6 (\href{https://doi.org/10.5281/zenodo.22721669}{doi:10.5281/zenodo.22721669}).
\end{fact}

\begin{fact}[The determinant identity, and the blow-up expansion]\label{c:sheaf.determinant-identity-blow-up-expansion}
Off the wall within the collar, $\det h = \det G \cdot \bigl(-(1 + 2\mu\nu\beta^2 s)\bigr)$ exactly, and $s = 1/\lambda + O(1)$ diverges with the sign of $\lambda$ as the wall is approached.
\emph{Source:} T4, Lem 4.3 (\href{https://doi.org/10.5281/zenodo.22721669}{doi:10.5281/zenodo.22721669}).
\end{fact}

\begin{fact}[Imported: classification values]\label{c:sheaf.imported-classification-values}
The real Clifford algebras $\Cl(1,3)$, $\Cl(0,4)$, $\Cl(0,3)$, $\Cl(2,4)$ and $\Cl(1,2)$ are isomorphic to $\Mat_2(\HH)$, $\Mat_2(\HH)$, $\HH \oplus \HH$, $\Mat_4(\HH)$ and $\Mat_2(\C)$ respectively.
\emph{Source:} T4, Fact 2.9 (\href{https://doi.org/10.5281/zenodo.22721669}{doi:10.5281/zenodo.22721669}).
\end{fact}

\begin{fact}[Imported: real line bundles and $w_1$]\label{c:sheaf.imported-real-line-bundles}
A real line bundle is trivial exactly when $w_1 = 0$, $w_1$ is additive so the square of any line bundle is trivial, and a hypersurface is two-sided exactly when its normal bundle is trivial.
\emph{Source:} J. W. Milnor and J. D. Stasheff, Characteristic Classes, Annals of Mathematics Studies 76, Princeton University Press (1974)., \S 4.
\end{fact}

\begin{fact}[The invariant form of (T3)]\label{c:sheaf.invariant-form-t3}
Given a smooth wall with a smooth corank-one radical line, the differential of $G(k,k)$ along the wall defines a canonical bundle map $\delta : \shL^{\otimes 2} \to N^*\Sig$, and transversality is exactly the demand that it be an isomorphism.
\emph{Source:} T4, Lem 2.2 (\href{https://doi.org/10.5281/zenodo.22721669}{doi:10.5281/zenodo.22721669}).
\end{fact}

\begin{fact}[Resolution, ideal-theoretically]\label{c:sheaf.resolution-ideal-theoretically}
In any nondegenerate quadratic space containing a wall fibre with its form, the radical line generates the unit ideal of the Clifford algebra.
\emph{Source:} T4, Prop 3.4 (\href{https://doi.org/10.5281/zenodo.22721669}{doi:10.5281/zenodo.22721669}).
\end{fact}

\begin{fact}[Tame wall induction]\label{c:sheaf.tame-wall-induction}
Assuming a compact tame wall and a good coupling, the shadow wall is a smooth closed hypersurface along which the ambient form is again a tame corank-one degeneration, so the completion iterates.
\emph{Source:} T4, Lem 5.3 (\href{https://doi.org/10.5281/zenodo.22721669}{doi:10.5281/zenodo.22721669}).
\end{fact}

\begin{fact}[The Clifford algebra exists and is unique]\label{c:std-algebra.existence-uniqueness}
$\Cl(V, Q) = T(V)/\langle v \otimes v - Q(v) \rangle$ exists and is unique up to unique isomorphism; the wedge-plus-contraction operators on $\Lambda V$ give a concrete representation.
\emph{Source:} the algebra core of the Transport Lab, proved there with an exact check; standard, see C. Chevalley, The Algebraic Theory of Spinors, Columbia University Press (1954), Ch. II (the Clifford algebra of a quadratic space).
\end{fact}

\begin{fact}[The radical ideal of a degenerate Clifford algebra and its quotient]\label{c:std-algebra.radical-ideal-and-quotient}
The radical generates the ideal $J = \Cl\operatorname{rad}V$, spanned by the blades with a radical factor, with $J^{r} \ne 0 = J^{r+1}$; $\Cl(V)/J \cong \Cl(V/\operatorname{rad}V)$, and $J$ is the largest nilpotent ideal.
\emph{Source:} the algebra core of the Transport Lab, proved there with an exact check; standard, see P. Lounesto, Clifford Algebras and Spinors, 2nd ed., London Mathematical Society Lecture Note Series 286, Cambridge University Press (2001), Ch. 14 (definitions of the Clifford algebra; degenerate quadratic forms).
\end{fact}

\begin{fact}[Orthogonal bases and the law of inertia]\label{c:std-algebra.sylvester-inertia}
Every real quadratic space has an orthogonal basis squaring to $1, -1, 0$, and the counts $(p, q, r)$ are invariants that classify it up to isometry.
\emph{Source:} T. Y. Lam, Introduction to Quadratic Forms over Fields, AMS (2005)., the diagonalisation of forms and Sylvester's law of inertia (real case).
\end{fact}

\begin{fact}[The product of two vectors]\label{c:std-algebra.vector-squares}
$v^{2} = Q(v)$ and $uv + vu = 2B(u, v)$, so $uv = B(u, v) + \tfrac12(uv - vu)$.
\emph{Source:} the algebra core of the Transport Lab, proved there with an exact check; standard, see C. Doran and A. Lasenby, Geometric Algebra for Physicists, Cambridge University Press (2003), Ch. 4 (foundations of geometric algebra).
\end{fact}

\section*{Conclusion}
The thread runs from an algebra that keeps its size but changes its structure on a wall, through a repair that gives the dying direction a partner, to shadow walls the repair can push away but never remove.

Everything rests on the tameness postulate in What we assume: the degeneracy forms a smooth closed wall, exactly one direction loses its length at each of its points, these directions vary smoothly, and the length vanishes at a linear rate. The one-sided repair also needs a positive direction on the physical side and a radical line that can be oriented consistently along the wall.

The endless sequence of walls further needs a compact wall and good couplings at every stage.

\end{document}